% 第九部：非核心情形與結論（第柒部分、結論，1401–1406）。出處：英文底稿；W06_1019_導讀問答.md「孤立而隔絕的少數」一則（已查核）。
\chapter{非核心情形與結論}

\qaentry{第九部之一}{假設不成立時論證的效力}

\begin{qaquestion}
四項假設不成立時，本文的論證如何？違憲審查是否因此就成立？
\end{qaquestion}

作者先答覆「本文的假設不切實際」之說：第三項假設相當容易滿足，因為違憲審查的支持者幾乎都預設該社會有一部與公民見解有實質關聯的權利法案；第一項假設並不要求完美\pg{1401--1402}。接著寫：

\begin{quote}
In cases in which the assumptions fail, the argument against judicial review presented in this Essay does not go through. As I emphasized in Part II, my argument is a conditional one. However, it does not follow that judicial review of legislation is defensible whenever the assumptions fail.\pg{1402}
\end{quote}

理由仍是比較性的：可能有不以這些假設為條件的其他反對論；違憲審查也可能無助於改善特定的處境；一個社會的司法若與其立法同樣腐敗、同樣帶有偏見，設立違憲審查並不適當\pg{1402}。

第一項假設不成立時，作者提出兩個問題：能否改善立法機關的處境？應否把重要權利問題的終局決定權交給法院\pg{1402--1403}？兩者並非全然獨立：法院作為後盾，可能使議員更難養成負責任的文化，作者自承 How far this is true, I don't know\pg{1403}。

\qaentry{第九部之二}{孤立而隔絕的少數}

\begin{qaquestion}
第柒部分為何以 Carolene Products 案腳註四收尾？作者如何界定 discrete 與 insular，又加了哪些限制？
\end{qaquestion}

1938 年 United States v.\ Carolene Products Co.\ 案腳註四寫：

\begin{quote}
[P]rejudice against discrete and insular minorities may be a special condition, which tends seriously to curtail the operation of those political processes ordinarily to be relied upon to protect minorities \ldots\pg{1403}
\end{quote}

作者稱這是描述違憲審查論據有若干說服力之非核心情形的極佳方式\pg{1403}，接著界定兩個形容詞\pg{1404}。

\begin{longtable}{@{}L{0.14\linewidth}L{0.42\linewidth}L{0.36\linewidth}@{}}
\toprule
用語 & 作者的界定 & 與第陸部分的關係 \\
\midrule
\endhead
discrete（孤立） & 獨立於政治決定而存在的少數，即議題上的少數 & 議題上的少數（1397） \\
insular（隔絕） & 成員與社群其他人少有共同利益，無從組成聯盟促進自身利益 & 兩種少數的重合（1398） \\
\bottomrule
\end{longtable}

對 Stone 所提的偏見要件，作者認為普遍的偏見與第三、四項假設不相容；重點在把各種敵意與對權利的合理分歧區分開來\pg{1404}。他另加三點限制：

\begin{enumerate}
\item 並非每個少數都適格，決定上的少數、議題上的少數、長期在表決中落敗者都不當然適格\pg{1403--1404}。
\item 要看司法上的多數是否同樣帶有偏見。支持違憲審查的論點假定同情少數權利的人集中於政治菁英，菁英中的法官較議員不受民意報復；同情若較多存在於一般人之間，選舉反而是把這份支持送進立法機關的管道\pg{1404--1405}。
\item 長期效果。美國 1950 與 1960 年代，法院攻擊種族隔離，最後起決定作用的是國會的民權立法，主要的分野在聯邦與州之間\pg{1405}。
\end{enumerate}

作者的結論是，不應把腳註四當作撬出一般違憲審查制度的槓桿\pg{1405--1406}。

\qaentry{第九部之三}{結論的兩項指控}

\begin{qaquestion}
結論把主張限定在什麼範圍？對違憲審查提出哪兩項指控？對支持者有什麼要求？
\end{qaquestion}

主張的範圍：

\begin{quote}
I have tried to show why rights-based judicial review is inappropriate for reasonably democratic societies whose main problem is not that their legislative institutions are dysfunctional but that their members disagree about rights.\pg{1406}
\end{quote}

結論第 2 段寫，一般的立法程序能以負責而審議的方式處理權利分歧，法院的終局審查在此之上所增加的，只是一種相當侮辱人的剝奪參政權（disenfranchisement），以及以法條主義遮蔽分歧中的道德問題（legalistic obfuscation）\pg{1406}。第 3 段承認，在特殊的病態、失靈的立法機關、腐敗的政治文化、種族主義與其他根深柢固偏見的遺緒之下，這些代價或許值得暫時承擔；但支持者應坦白以此為理由，並帶著 a degree of humility and shame，不應把違憲審查當作尊重權利的典範向外推廣\pg{1406}。

\begin{readernote}
結論兩項指控與導論的兩條戰線\pg{1353}一一對應：法律技術的遮蔽對應第肆部分，剝奪參政權對應第伍部分。導讀收尾時可以把導論與結論並排，讓聽者看到全文首尾的對應。
\end{readernote}
